% AA-36 | Audyt historyczny VI: pełne dowody i przykłady
\chapter{Wielomian ogólny i jego grupa Galois}\label{aa:ch:generic-polynomial}
\section{Algebraiczna niezależność współczynników}
Dla ustalonego wielomianu nad $\mathbb Q$ grupa Galois zależy od
jego współczynników. Na przykład dwa nierozkładalne wielomiany stopnia
trzeciego mogą mieć grupy $S_3$ i $C_3$. Chcemy oddzielić własności
wspólne \emph{wszystkim formalnie niezależnym współczynnikom} od
szczególnych równań, które pojawiają się po podstawieniu konkretnych
liczb. W tym celu zamiast wybierać współczynniki w $k$, wprowadzamy
zmienne transcendentalne nad $k$.

\begin{definicja}\label{aa:def:generic-poly}
Niech $k$ będzie ciałem, $n\ge2$ liczbą całkowitą, a
$t_1,\ldots,t_n$ zmiennymi algebraicznie niezależnymi nad $k$;
znaczy to, że dla każdego niezerowego
$P\in k[T_1,\ldots,T_n]$ zachodzi $P(t_1,\ldots,t_n)\ne0$.
Przez $k(t_1,\ldots,t_n)$ rozumiemy ciało ułamków
$k[t_1,\ldots,t_n]$. Wielomian
\[
 F(X)=X^n-t_1X^{n-1}+t_2X^{n-2}-\cdots+(-1)^nt_n
 \in k(t_1,\ldots,t_n)[X]
\]
nazywamy \emph{ogólnym wielomianem monicznym stopnia $n$}.
Alternatywna konwencja znaków, z niezależnymi współczynnikami
przy kolejnych potęgach $X$, daje ten sam typ wielomianu: zmiana
$t_i\mapsto(-1)^it_i$ jest automorfizmem pierścienia wielomianów.
\end{definicja}

\begin{lemat}[Niezależność elementarnych funkcji symetrycznych]
\label{aa:lem:generic-independence}
Niech $x_1,\ldots,x_n$ będą algebraicznie niezależne nad ciałem $k$.
Dla $1\le j\le n$ oznaczmy
$e_j=\sum_{1\le i_1<\cdots<i_j\le n}x_{i_1}\cdots x_{i_j}$.
Wówczas $e_1,\ldots,e_n$ są algebraicznie niezależne nad $k$.
\end{lemat}
\begin{proof}
Rozważmy homomorfizm
$\varphi:k[T_1,\ldots,T_n]\to k[X_1,\ldots,X_n]$
określony wzorem $\varphi(T_j)=e_j(X_1,\ldots,X_n)$.
Podstawowe twierdzenie o wielomianach symetrycznych
(Twierdzenie~\ref{aa:tw:fundamental-symmetric}) stwierdza nie tylko,
że każdy wielomian symetryczny jest wielomianem w $e_j$, lecz także,
że taki zapis jest \emph{jednoznaczny}. Stąd $\ker\varphi=(0)$.
Ponadto odwzorowanie
$\psi:k[X_1,\ldots,X_n]\to k(x_1,\ldots,x_n)$
określone przez $X_i\mapsto x_i$ jest injektywne: gdyby
$P(x_1,\ldots,x_n)=0$ dla niezerowego $P$, otrzymalibyśmy relację
algebraiczną między $x_i$. Zatem także $\psi\circ\varphi$ jest
injektywne. Jest to dokładnie warunek niezależności algebraicznej
$e_1,\ldots,e_n$.
\end{proof}

\section{Obliczenie grupy Galois}
\begin{twierdzenie}[Grupa Galois wielomianu ogólnego]
\label{aa:tw:generic-sn}
Dla każdego ciała $k$ i każdej liczby całkowitej $n\ge2$ ogólny
wielomian moniczny $F\in k(t_1,\ldots,t_n)[X]$ jest nierozkładalny
i rozdzielczy, a jego grupa Galois nad $k(t_1,\ldots,t_n)$
jest izomorficzna z grupą symetryczną $S_n$.
\end{twierdzenie}
\begin{proof}
\emph{Krok 1: ciało zawierające formalne pierwiastki.}
Weźmy algebraicznie niezależne zmienne $x_1,\ldots,x_n$ nad $k$,
połóżmy $L=k(x_1,\ldots,x_n)$ i
$K=k(e_1(x_1,\ldots,x_n),\ldots,e_n(x_1,\ldots,x_n))$.
Lemat~\ref{aa:lem:generic-independence} daje izomorfizm
$k(t_1,\ldots,t_n)\cong K$, określony przez $t_j\mapsto e_j$.
Będziemy więc obliczać grupę Galois nad $K$.
Ze wzorów Viète'a wynika tożsamość
\[
 f(X)=\prod_{i=1}^n(X-x_i)
 =X^n-e_1X^{n-1}+\cdots+(-1)^ne_n\in K[X].
\]
Każdy $x_i$ jest pierwiastkiem $f$; ponieważ są to różne zmienne,
$x_i-x_j\ne0$ dla $i\ne j$, więc $f$ ma $n$ różnych pierwiastków.
Jest zatem rozdzielczy również w dodatniej charakterystyce.
Ponadto $L=K(x_1,\ldots,x_n)$: z $k\subseteq K$ mamy
$L=k(x_1,\ldots,x_n)\subseteq K(x_1,\ldots,x_n)$,
a z $K\subseteq L$ wynika odwrotne zawieranie.
Tym samym $L$ jest ciałem rozkładu $f$ nad $K$.

\emph{Krok 2: permutacje zmiennych rzeczywiście dają automorfizmy.}
Dla $\sigma\in S_n$ podstawienie $x_i\mapsto x_{\sigma(i)}$
indukuje automorfizm pierścienia $k[x_1,\ldots,x_n]$:
jego odwrotnością jest podstawienie wyznaczone przez $\sigma^{-1}$.
Ponieważ automorfizm przenosi niezerowe wielomiany na niezerowe,
rozszerza się jednoznacznie na ciało ułamków $L$ przez
$P/Q\mapsto\sigma(P)/\sigma(Q)$.
Wielomiany $e_j$ są symetryczne, więc $\sigma(e_j)=e_j$;
automorfizm ustala punktowo $K$.
Różne permutacje dają różne automorfizmy, bo dla pewnego $i$
mają różne wartości na $x_i$. Otrzymaliśmy więc co najmniej
$n!$ elementów $\operatorname{Gal}(L/K)$.

\emph{Krok 3: nie ma innych automorfizmów.}
Jeżeli $\tau$ jest automorfizmem $L$ ustalającym $K$, to
$f(\tau(x_i))=\tau(f(x_i))=0$.
Wobec rozdzielczości $f$ obraz $\tau(x_i)$ jest jednym z $x_j$.
Automorfizm jest injektywny, więc obrazy $n$ różnych pierwiastków
są parami różne. Stąd $\tau$ wyznacza permutację zbioru
$\{x_1,\ldots,x_n\}$; ponieważ ten zbiór generuje $L$ nad $K$,
permutacja całkowicie wyznacza $\tau$. Mamy więc co najwyżej $n!$
automorfizmów, a z Kroku 2 otrzymujemy
$\operatorname{Gal}(L/K)\cong S_n$.

\emph{Krok 4: nierozkładalność.}
Gdyby $f=gh$ z $g,h\in K[X]$ o dodatnich stopniach,
zbiory ich pierwiastków w $L$ byłyby niepuste i właściwe
(ponieważ $f$ jest rozdzielczy).
Dla automorfizmu $\tau\in\operatorname{Gal}(L/K)$ oraz pierwiastka
$a$ wielomianu $g$ mielibyśmy $g(\tau(a))=\tau(g(a))=0$.
Zbiór pierwiastków $g$ byłby więc niezmienniczy względem
całej grupy $S_n$. Jest to niemożliwe, gdyż $S_n$ działa
przechodnio na $\{x_1,\ldots,x_n\}$, a jedyne niezmiennicze
podzbiory zbioru będącego jedną orbitą to zbiór pusty i cały zbiór.
Zatem $f$ jest nierozkładalny.
\end{proof}

\begin{wniosek}[Ogólna niemożliwość rozwiązania przez pierwiastniki]
\label{aa:wn:generic-abel}
Jeżeli $\operatorname{char}k=0$ i $n\ge5$, to ogólny wielomian
stopnia $n$ nie jest rozwiązywalny przez pierwiastniki nad
$k(t_1,\ldots,t_n)$.
\end{wniosek}
\begin{proof}
Z Twierdzenia~\ref{aa:tw:generic-sn} jego grupa Galois to $S_n$.
Dla $n\ge5$ grupa $S_n$ zawiera podgrupę izomorficzną
z $A_5$: parzyste permutacje indeksów $1,\ldots,5$
przedłużamy, pozostawiając indeksy $6,\ldots,n$ nieruchome.
Wniosek~\ref{aa:wn:a5-not-solvable}, udowodniony na podstawie
prostoty i nieprzemienności $A_5$, stwierdza,
że $A_5$ nie jest rozwiązalna.
Każda podgrupa grupy rozwiązalnej jest rozwiązalna:
jej szereg pochodny zawiera się wyraz po wyrazie
w szeregu pochodnym całej grupy.
Zatem $S_n$ nie może być rozwiązalna.
Kryterium z Rozdziału~\ref{aa:ch:radicals} (w charakterystyce zero)
mówi, że rozwiązywalność wielomianu rozdzielczego przez pierwiastniki
pociąga rozwiązalność jego grupy Galois.
Stosując przeciwstawienie, otrzymujemy tezę.
\end{proof}

\begin{wniosek}[Ciało niezmienników i stopień rozszerzenia]
\label{aa:wn:generic-invariants}
Niech $k$ będzie dowolnym ciałem, $n\ge2$ i niech
$x_1,\ldots,x_n$ będą algebraicznie niezależnymi elementami
nad $k$. Oznaczmy $L=k(x_1,\ldots,x_n)$ i
$K=k(e_1,\ldots,e_n)$, gdzie $e_j$ jest $j$-tą elementarną
funkcją symetryczną zmiennych $x_1,\ldots,x_n$.
Wówczas
\[
 k(x_1,\ldots,x_n)^{S_n}=k(e_1,\ldots,e_n),
 \qquad [L:K]=n!,
\]
gdzie $S_n$ działa przez permutacje zmiennych. Innymi słowy,
każda funkcja wymierna $x_1,\ldots,x_n$ niezmiennicza względem
wszystkich ich permutacji jest funkcją wymierną elementarnych
funkcji symetrycznych.
\end{wniosek}
\begin{proof}
W dowodzie Twierdzenia~\ref{aa:tw:generic-sn} wykazano, że
$L/K$ jest skończonym rozszerzeniem Galois i że jego grupa
Galois jest dokładnie grupą wszystkich permutacji zmiennych,
czyli $S_n$. Z podstawowego twierdzenia teorii Galois
(Twierdzenie~\ref{aa:tw:fundamental-galois}) wynika
$L^{\operatorname{Gal}(L/K)}=K$. Ponieważ działanie
grupy $S_n$ pokrywa się z działaniem całej grupy Galois,
pierwsza równość jest udowodniona.
Rozszerzenie $L/K$ jest skończone i galoisowskie,
więc stopień rozszerzenia równa się rzędowi grupy Galois:
$[L:K]=|S_n|=n!$.
\end{proof}

\begin{przyklad}[Dwie funkcje wymierne niezmiennicze względem permutacji]
\label{aa:prz:generic-rational-invariants}
Niech $x_1,x_2,x_3$ będą algebraicznie niezależne nad ciałem
$k$, a $e_1,e_2,e_3$ odpowiadającymi im elementarnymi
funkcjami symetrycznymi. Rozwijając kwadrat sumy, otrzymujemy
\[
 x_1^2+x_2^2+x_3^2=(x_1+x_2+x_3)^2
 -2(x_1x_2+x_1x_3+x_2x_3)=e_1^2-2e_2.
\]
Tożsamość obowiązuje również w charakterystyce $2$,
gdzie jej ostatni składnik znika. Ponieważ algebraiczna
niezależność gwarantuje $x_1x_2x_3=e_3\ne0$, można też
przeprowadzić rachunek
\[
 \frac1{x_1}+\frac1{x_2}+\frac1{x_3}
 =\frac{x_2x_3+x_1x_3+x_1x_2}{x_1x_2x_3}
 =\frac{e_2}{e_3}.
\]
Obie funkcje pozostają niezmienione po dowolnej permutacji
trzech zmiennych; uzyskane wzory jawnie przedstawiają je
jako funkcje wymierne $e_1,e_2,e_3$.
\end{przyklad}

\begin{przyklad}[Kwadratowe ciało pośrednie dla wielomianu ogólnego stopnia trzeciego]
\label{aa:prz:generic-discriminant}
Niech $\operatorname{char}k\ne2$ i $n=3$.
W ciele $L=k(x_1,x_2,x_3)$ połóżmy
$\Delta=(x_1-x_2)(x_1-x_3)(x_2-x_3)$.
Dla każdej permutacji $\sigma\in S_3$ zachodzi
$\sigma(\Delta)=\operatorname{sgn}(\sigma)\Delta$.
W szczególności $\Delta$ jest ustalany przez wszystkie
elementy $A_3$, więc należy do $L^{A_3}$.
Jego kwadrat $D=\Delta^2$ jest niezmienniczy względem
całej grupy $S_3$ i dlatego należy do
$K=k(e_1,e_2,e_3)$. Rozwijając iloczyn różnic,
otrzymujemy jawnie
\[
 D=e_1^2e_2^2-4e_2^3-4e_1^3e_3-27e_3^2+18e_1e_2e_3.
\]
Jednak $\Delta\notin K$:
transpozycja zamienia go na $-\Delta\ne\Delta$.
Ponieważ $\operatorname{char}k\ne2$, jedynymi pierwiastkami
$T^2-D$ w $L$ są $\Delta$ i $-\Delta$.
Żaden z nich nie należy do $K$, bo transpozycja zmienia
znak niezerowego $\Delta$. Wielomian $T^2-D\in K[T]$
nie ma więc pierwiastka w $K$ i, jako wielomian stopnia dwa,
jest nierozkładalny.
Otrzymujemy $[K(\Delta):K]=2$. Podstawowe twierdzenie
teorii Galois daje $[L^{A_3}:K]=[S_3:A_3]=2$, więc
$L^{A_3}=K(\Delta)$. Jest to przykład nietrywialnego
ciała pośredniego, wyznaczonego przez parzyste permutacje.
\end{przyklad}

\section{Przykłady i granice twierdzenia}
\begin{przyklad}[Dwa nierozkładalne wielomiany sześcienne]
\label{aa:prz:generic-cubic}
Dla algebraicznie niezależnych $a,b,c$ nad $\mathbb Q$
wielomian $X^3+aX^2+bX+c$ ma grupę $S_3$ nad $\mathbb Q(a,b,c)$:
w Definicji~\ref{aa:def:generic-poly} wystarczy zastąpić
$t_1=-a$, $t_2=b$, $t_3=-c$.
Po podstawieniu konkretnych współczynników sytuacja się zmienia.
Wielomian $g(X)=X^3-2$ jest nierozkładalny nad $\mathbb Q$
(kryterium Eisensteina dla $2$), a jego wyróżnik wynosi
$-4\cdot0^3-27\cdot(-2)^2=-108$.
Ponieważ $-108$ nie jest kwadratem w $\mathbb Q$,
grupa Galois działająca przechodnio na trzech pierwiastkach
nie zawiera się w $A_3$, więc jest równa $S_3$.
Natomiast $h(X)=X^3-3X+1$ nie ma pierwiastków wymiernych:
jedynymi kandydatami według twierdzenia o pierwiastkach wymiernych
są $\pm1$, a $h(1)=-1$ i $h(-1)=3$.
Stąd $h$ jest nierozkładalny, a jego wyróżnik jest równy
$-4(-3)^3-27\cdot1^2=81=9^2$.
W charakterystyce zero kryterium wyróżnika daje więc podgrupę
$A_3$; z przechodniości wynika, że jest ona całym
$A_3\cong C_3$. Oba wielomiany mają stopień trzy, lecz różne
grupy Galois.
\end{przyklad}

\begin{przyklad}[Specjalizacja może zmniejszyć grupę]
\label{aa:prz:generic-specialization}
Wielomian ogólny stopnia czwartego ma grupę $S_4$.
Dla specjalizacji do $X^4-2$ otrzymujemy natomiast grupę
izomorficzną z grupą dihedralną $D_4$ rzędu $8$:
po wprowadzeniu $\alpha=\sqrt[4]{2}$ ciałem rozkładu jest
$\mathbb Q(\alpha,i)$, a jej osiem automorfizmów opisano
w Przykładzie~\ref{aa:prz:quartic-v4}.
Nie jest to sprzeczność: specjalizacja narzuca zależności między
współczynnikami, więc przestają one być algebraicznie niezależne.
\end{przyklad}

\begin{przyklad}[Dlaczego nie pomijamy rozdzielczości]
\label{aa:prz:generic-inseparable-specialization}
Niech $k=\mathbb F_p(t)$, gdzie $p$ jest liczbą pierwszą,
a $t$ elementem transcendentalnym nad $\mathbb F_p$.
Wielomian $q(X)=X^p-t\in k[X]$ jest nierozkładalny.
Istotnie, rozpatrując go najpierw w $\mathbb F_p[t][X]$,
widzimy, że wszystkie współczynniki oprócz wiodącego są
podzielne przez pierwszy element $t$, a wyraz stały $-t$
nie jest podzielny przez $t^2$. Kryterium Eisensteina
w pierścieniu $\mathbb F_p[t]$ i lemat Gaussa dają
nierozkładalność w jego ciele ułamków $k=\mathbb F_p(t)$.
Jeżeli $\alpha$ jest jego pierwiastkiem w domknięciu
algebraicznym $k$, to z $\alpha^p=t$ wynika
$q(X)=(X-\alpha)^p$, ponieważ w charakterystyce $p$
zachodzi $(X-\alpha)^p=X^p-\alpha^p$.
Stąd $q'(X)=0$ i w ciele rozkładu występuje tylko jeden
różny pierwiastek $\alpha$, mimo że $q$ jest nierozkładalny.
Jego grupa automorfizmów nad $k$
jest trywialna: każdy $k$-automorfizm ciała rozkładu
$k(\alpha)$ musi wysłać $\alpha$ na pierwiastek $q$,
a ten pierwiastek jest jedyny, zatem automorfizm ustala
$\alpha$ i całe $k(\alpha)$. Nie otrzymujemy więc $S_p$.
To przykład \emph{specjalizacji}, a nie kontrprzykład do
Twierdzenia~\ref{aa:tw:generic-sn}: w modelu użytym w jego dowodzie
pierwiastki są różnymi zmiennymi algebraicznie niezależnymi nad $k$,
natomiast przy $n=p$ podstawienie $t_1=\cdots=t_{p-1}=0$,
$t_p=t$ daje właśnie $X^p-t$, który nie jest rozdzielczy.
\end{przyklad}

\begin{uwaga}\label{aa:uw:generic-characteristic}
Twierdzenie~\ref{aa:tw:generic-sn} działa nad ciałem dowolnej
charakterystyki, ponieważ konstrukcja permutacji zmiennych nie
wymaga dzielenia przez $n!$. Natomiast Wniosek
\ref{aa:wn:generic-abel} opiera się na kryterium rozwiązywalności
przez pierwiastniki przyjętym w Rozdziale~\ref{aa:ch:radicals},
stąd jawne założenie charakterystyki zero.
\end{uwaga}
